\documentclass[12pt,a4paper]{amsart}

\usepackage{amsmath,amssymb,amsthm,mathtools,enumitem,booktabs,array,hyperref}
\usepackage[utf8]{inputenc}
\usepackage[T1]{fontenc}
\usepackage[ngerman]{babel}
\usepackage{geometry}
\geometry{margin=2.5cm}

% -------------------------------------------------------
% Theorem-Umgebungen (Englisch)
% -------------------------------------------------------
\newtheorem{theorem}{Theorem}[section]
\newtheorem{lemma}[theorem]{Lemma}
\newtheorem{corollary}[theorem]{Korollar}
\newtheorem{proposition}[theorem]{Proposition}
\theoremstyle{definition}
\newtheorem{definition}[theorem]{Definition}
\newtheorem{example}[theorem]{Beispiel}
\newtheorem{remark}[theorem]{Remark}
\newtheorem{conjecture}[theorem]{Conjecture}

% Deutschsprachige Theorem-Umgebungen
\newtheorem{satz}[theorem]{Satz}
\newtheorem{vermutung}[theorem]{Vermutung}
\newtheorem{korollar}[theorem]{Korollar}
\newtheorem{bemerkung}[theorem]{Bemerkung}
\newtheorem{definition_de}[theorem]{Definition}

% -------------------------------------------------------
% Mathematische Abkürzungen
% -------------------------------------------------------
\newcommand{\N}{\mathbb{N}}
\newcommand{\Z}{\mathbb{Z}}
\newcommand{\Q}{\mathbb{Q}}
\newcommand{\R}{\mathbb{R}}
\newcommand{\C}{\mathbb{C}}
\newcommand{\Ac}{\mathcal{A}}
\newcommand{\Bc}{\mathcal{B}}
\newcommand{\At}{\widehat{\mathcal{A}}}
\newcommand{\dr}{\,\mathrm{d}}
\newcommand{\KZ}{\mathrm{KZ}}
\DeclareMathOperator{\tr}{Spur}
\DeclareMathOperator{\End}{End}

% -------------------------------------------------------
\begin{document}
% -------------------------------------------------------

\title{Das Kontsevich-Integral und Vassiliev-Invarianten:\\
       Klassifikation von Knoten durch Sehnendiagramme}

\author{Michael Fuhrmann}
\address{Specialist Mathematics Project, Build 170}
\email{specialist-maths@localhost}
\date{\today}

\begin{abstract}
Vassiliev-Invarianten (Invarianten endlichen Typs), eingeführt von Vassiliev
im Jahr 1990 und kombinatorisch neu formuliert von Birman und Lin, gehören zu den
mächtigsten und strukturreichsten Invarianten der nieder-dimensionalen Topologie.
Das Kontsevich-Integral $Z(K)$, konstruiert von Maxim Kontsevich 1993, ist eine
universelle Vassiliev-Invariante: Es ist eine formale Potenzreihe in der graduierten
Algebra $\mathcal{A}$ der Sehnendiagramme modulo der $4T$- und $1T$-Relationen,
durch die jede Vassiliev-Invariante faktoriert.  Wir entwickeln den vollständigen
Rahmen: Sehnendiagramme und $4T/1T$-Relationen, Gewichtssysteme und ihre
Lie-algebraische Konstruktion, das Kontsevich-Integral als iteriertes Chen-Integral,
der Drinfeld-Assoziator $\Phi_{\mathrm{KZ}}$ und seine Rolle bei der Regularisierung,
das farbige Jones-Polynom als Spezialisierung, die offene Vermutung über die
Vollständigkeit von $Z$, der Wheeling-Satz (Bar-Natan--Le--Thurston) und die
Le--Murakami--Ohtsuki-Invariante.  Alle unbewiesenen Aussagen werden sorgfältig
als Vermutungen gekennzeichnet.
\end{abstract}

\maketitle

\tableofcontents

% -------------------------------------------------------
\section{Einleitung und historischer Überblick}
% -------------------------------------------------------

Die Klassifikation von Knoten --- die topologische Klasse von Einbettungen
$K: S^1 \hookrightarrow S^3$ bis auf ambiante Isotopie --- ist ein zentrales
Problem der nieder-dimensionalen Topologie.  Klassische Invarianten wie das
Alexander-Polynom (1928) und das Jones-Polynom (1984) unterscheiden viele
Knoten, sind aber keine vollständigen Invarianten.

Die Entdeckung des Jones-Polynoms~\cite{jones1985} löste eine Revolution aus:
Witten~\cite{witten1989} interpretierte es über dreidimensionale
Chern-Simons-Eichtheorie mit Eichgruppe $\mathrm{SU}(2)$, was andeutete,
dass eine ``vollständige'' Knoteninvariante aus der Quantenfeldtheorie stammen
könnte.  Gleichzeitig führte Vassiliev~\cite{vassiliev1990} Invarianten
endlichen Typs über die Topologie von Diskriminantenräumen (Räumen singulärer
Knoten) ein.

Kontsevich~\cite{kontsevich1993} vereinigte diese Ansätze durch die Konstruktion
des \emph{Kontsevich-Integrals}: einer einzigen formal-potenzreihen-wertigen
Invariante, die alle Vassiliev-Invarianten als Koeffizienten ihrer Entwicklung
enthält.  Kontsevichs Konstruktion stützt sich auf die Theorie der iterierten
Chen-Integrale, den Drinfeld-Assoziator aus der Quantengruppen-Theorie und
die kombinatorischen Strukturen der Sehnendiagramme.

\subsection*{Gliederung}

Abschnitt~\ref{sec:vassiliev} definiert Vassiliev-Invarianten über die
Vassiliev-Auflösung und singuläre Knoten.
Abschnitt~\ref{sec:sehnen} führt Sehnendiagramme, die $4T$- und $1T$-Relationen
sowie die Algebra $\Ac$ ein.
Abschnitt~\ref{sec:gewicht} entwickelt Gewichtssysteme und ihre Lie-algebraische
Konstruktion.
Abschnitt~\ref{sec:KI} konstruiert das Kontsevich-Integral via iterierter Integrale.
Abschnitt~\ref{sec:konvergenz} behandelt Konvergenz und den Drinfeld-Assoziator.
Abschnitt~\ref{sec:assoziator} behandelt den Knizhnik-Zamolodchikov-Assoziator
$\Phi_{\mathrm{KZ}}$.
Abschnitt~\ref{sec:jones} stellt den Bezug zum farbigen Jones-Polynom her.
Abschnitt~\ref{sec:universalitaet} beweist die Universalität des Kontsevich-Integrals.
Abschnitt~\ref{sec:vollstaendigkeit} diskutiert die Vollständigkeitsvermutung.
Abschnitt~\ref{sec:wheeling} präsentiert den Wheeling-Satz.
Abschnitt~\ref{sec:LMO} führt die Le-Murakami-Ohtsuki-Invariante ein.
Abschnitt~\ref{sec:chern} verbindet die Theorie mit der Chern-Simons-Feldtheorie.
Abschnitt~\ref{sec:offen} listet offene Probleme.

% -------------------------------------------------------
\section{Vassiliev-Invarianten (Invarianten endlichen Typs)}\label{sec:vassiliev}
% -------------------------------------------------------

\begin{definition_de}[Knoten-Invariante]
Ein \emph{Knoten} ist eine ambiante Isotopie-Klasse glatter Einbettungen
$S^1 \to S^3$ (oder $\R^3$).  Eine \emph{Knoten-Invariante} ist eine
Abbildung $v: \{\text{Knoten}\} \to \mathbf{k}$ (wobei $\mathbf{k}$ ein
kommutativer Ring ist, typischerweise $\Q$ oder $\C$), die unter ambianter
Isotopie invariant ist.
\end{definition_de}

\begin{definition_de}[Singuläre Knoten und die Vassiliev-Scheinrelation]
\label{def:vassiliev}
Ein \emph{singulärer Knoten mit $n$ Doppelpunkten} ist eine glatte Immersion
$K: S^1 \looparrowright S^3$ mit genau $n$ transversalen Selbstschnitten.
Erweitere jede Knoten-Invariante $v$ auf singuläre Knoten durch die
\emph{Vassiliev-Relation} (Differenzenformel): An jedem markierten
Doppelpunkt $\times$ setze
\[
  v(K_\times) := v(K_+) - v(K_-),
\]
wobei $K_+$ (bzw.\ $K_-$) durch positive (bzw.\ negative) Auflösung des
Doppelpunkts entsteht.  Dann heißt $v$ eine \emph{Vassiliev-Invariante
vom Typ $\leq n$} (oder \emph{Invariante endlichen Typs der Ordnung $n$}),
wenn $v$ auf allen singulären Knoten mit mehr als $n$ Doppelpunkten
verschwindet:
\[
  v(K_\times) = 0 \quad \text{für alle singulären Knoten mit } > n \text{ Doppelpunkten.}
\]
\end{definition_de}

\begin{example}[Vassiliev-Invarianten kleiner Ordnung]
\begin{enumerate}[label=(\alph*)]
  \item Jede Invariante vom Typ $0$ ist konstant.
  \item Jede Invariante vom Typ $1$ ist konstant (es gibt keine nicht-trivialen
        Invarianten vom Typ 1).
  \item Die erste nicht-triviale Vassiliev-Invariante hat Typ $2$; es ist der
        Koeffizient $a_2$ im Conway-Polynom $\nabla(K)(z) = \sum a_{2i} z^{2i}$.
        Äquivalent dazu ist die Casson-Invariante $\lambda(K)$ eine
        Vassiliev-Invariante vom Typ 2.
  \item Das farbige Jones-Polynom $J_N(K;e^h) = \sum_n v_n^{(N)}(K) h^n$
        liefert Vassiliev-Invarianten $v_n^{(N)}$ vom Typ $\leq n$ für jeden
        ganzzahligen Parameter $N \geq 1$.
\end{enumerate}
\end{example}

\begin{satz}[Birman-Lin 1993, Taylor-Entwicklung von Knoten-Invarianten]
\label{satz:birmanlin}
Vassiliev-Invarianten der Ordnung $\leq n$ sind genau diejenigen
Knoten-Invarianten, deren ``$n$-te Ableitung am Unknoten'' wohldefiniert
ist.  Sie bilden einen filtrierten Vektorraum
$\mathcal{V}_0 \subseteq \mathcal{V}_1 \subseteq \mathcal{V}_2 \subseteq \cdots$
mit $\bigcup_n \mathcal{V}_n = \mathcal{V}$.
\end{satz}

% -------------------------------------------------------
\section{Sehnendiagramme und die Algebra $\mathcal{A}$}\label{sec:sehnen}
% -------------------------------------------------------

\begin{definition_de}[Sehnendiagramm]
Ein \emph{Sehnendiagramm der Ordnung $n$} ist ein orientierter Kreis
(der das Definitionsgebiet $S^1$ darstellt) mit $n$ ungeordneten Paaren
verschiedener Punkte, die durch \emph{Sehnen} (Bögen im Innern der Scheibe)
verbunden sind.  Zwei Sehnendiagramme sind gleich, wenn sie durch einen
orientierungserhaltenden Diffeomorphismus des Kreises zusammenhängen.
Sei $\Ac_n$ der $\Q$-Vektorraum, aufgespannt von Sehnendiagrammen der
Ordnung $n$.
\end{definition_de}

\begin{definition_de}[$4T$-Relation]
Die \emph{Vier-Term-Relation ($4T$)} in $\Ac_n$ lautet:
\[
  D_1 - D_2 + D_3 - D_4 = 0,
\]
wobei $D_1, D_2, D_3, D_4$ vier Sehnendiagramme sind, die außerhalb eines
festen lokalen Bereichs übereinstimmen und in diesem Bereich auf die vier
Arten verbunden sind, zwei Stränge durch eine zusätzliche Sehne zu verbinden.
Für zwei Sehnen mit gemeinsamem Endpunkt:
\[
  [ab] - [ac] + [bc] - [bd] = 0
\]
in einem Bereich mit Strängen $a, b, c, d$ von links nach rechts.
\end{definition_de}

\begin{definition_de}[$1T$-Relation (Gestellunabhängigkeit)]
Die \emph{Ein-Term-Relation ($1T$)} besagt, dass jedes Sehnendiagramm,
das eine \emph{isolierte Sehne} enthält (eine Sehne, deren beide Endpunkte
benachbart auf dem Kreis sind, ohne andere Sehnenendpunkte dazwischen),
gleich null ist.
\end{definition_de}

\begin{definition_de}[Sehnendiagramm-Algebra $\mathcal{A}$]
Die \emph{Sehnendiagramm-Algebra} ist die graduierte $\Q$-Algebra:
\[
  \Ac = \bigoplus_{n=0}^\infty \Ac_n, \quad
  \Ac_n = \frac{\Q[\text{Sehnendiagramme der Ordnung } n]}{4T + 1T \text{-Relationen}}.
\]
Die Multiplikation ist durch die verbundene Summe von Sehnendiagrammen gegeben:
$D \cdot D' = D \# D'$ (Verkettung entlang des Kreises).
\end{definition_de}

\begin{satz}[Dimensionen von $\mathcal{A}_n$]
Die ersten Dimensionen sind:
\[
  \dim \Ac_0 = 1, \quad \dim \Ac_1 = 0, \quad \dim \Ac_2 = 1,
  \quad \dim \Ac_3 = 1, \quad \dim \Ac_4 = 3, \quad \dim \Ac_5 = 4.
\]
Die erzeugende Funktion $\sum_n (\dim \Ac_n) t^n$ ist noch nicht in
geschlossener Form verstanden.
\end{satz}

\begin{bemerkung}[$STU$-Relation und Jacobi-Diagramme]
Eine wichtige Erweiterung ersetzt Sehnendiagramme durch \emph{Jacobi-Diagramme}
(auch chinesische Zeichendiagramme oder $3$-Graphen genannt): Diagramme mit
einwertigen Knoten auf einem Kreis und dreiwertigen Knoten im Inneren.
Diese erfüllen die \emph{STU-Relation} $S = T - U$, und die resultierende
Algebra $\mathcal{B}$ ist zu $\Ac$ isomorph (der Poincaré-Birkhoff-Witt-Isomorphismus
für Jacobi-Diagramme).
\end{bemerkung}

% -------------------------------------------------------
\section{Gewichtssysteme und Lie-Algebren}\label{sec:gewicht}
% -------------------------------------------------------

\begin{definition_de}[Gewichtssystem]
Ein \emph{Gewichtssystem der Ordnung $n$} ist eine lineare Abbildung
$W: \Ac_n \to \mathbf{k}$ (die die $4T$- und $1T$-Relationen erfüllt,
d.h.\ $W \in \Ac_n^*$).  Ein Gewichtssystem aller Ordnungen ist eine
Folge $\{W_n\}$ mit $W_n \in \Ac_n^*$ für jedes $n$.
\end{definition_de}

\begin{satz}[Vassiliev-Invarianten $\leftrightarrow$ Gewichtssysteme,
Vassiliev-Kontsevich]\label{satz:vs_ws}
Es gibt einen kanonischen Isomorphismus gefilterter Vektorräume:
\[
  \frac{\mathcal{V}_n}{\mathcal{V}_{n-1}} \xrightarrow{\;\sim\;} \Ac_n^*.
\]
Jede Vassiliev-Invariante $v \in \mathcal{V}_n$ bestimmt ein wohldefiniertes
``Symbol'' in $\Ac_n^*$, und jedes Element von $\Ac_n^*$ entsteht auf
diese Weise.
\end{satz}

\begin{satz}[Lie-algebraische Gewichtssysteme, Bar-Natan 1995]
\label{satz:lie_ws}
Sei $\mathfrak{g}$ eine metriesierte endlich-dimensionale Lie-Algebra
(mit invarianter, nicht-entarteter Bilinearform $(\cdot,\cdot)$), und
sei $R$ eine endlich-dimensionale Darstellung von $\mathfrak{g}$.
Definiere das Gewichtssystem
$W_{\mathfrak{g},R}: \Ac_n \to \mathbf{k}$ durch:
\[
  W_{\mathfrak{g},R}(D) = \mathrm{Spur}_R\!\left(\prod_{\text{Sehnen } (i,j) \text{ von } D}
  \Omega_{ij}\right),
\]
wobei $\Omega = \sum_a e_a \otimes e^a$ das Casimir-Element ist.
Dann ist $W_{\mathfrak{g},R}$ ein wohldefiniertes Gewichtssystem.
\end{satz}

\begin{example}[Gewichtssysteme aus $\mathfrak{sl}_2$]
Für $\mathfrak{g} = \mathfrak{sl}_2(\C)$ mit der 2-dimensionalen
(Standard-) Darstellung liefert das Gewichtssystem die Koeffizienten
des farbigen Jones-Polynoms $J_2$.  Für die $n$-dimensionale Darstellung
erhält man $J_n$ (das $n$-farbige Jones-Polynom).
\end{example}

\begin{vermutung}[Alle Gewichtssysteme sind Lie-algebraisch]\label{verm:alle_lie}
Jedes Gewichtssystem stammt von einer endlich-dimensionalen metriseierten
Lie-Algebra und einer ihrer Darstellungen.  Äquivalent: Alle Vassiliev-Invarianten
kommen von Quantengruppen.
\end{vermutung}

% -------------------------------------------------------
\section{Das Kontsevich-Integral}\label{sec:KI}
% -------------------------------------------------------

Sei $K \subset \R^2 \times \R$ ein orientierter Knoten in ``allgemeiner Lage''
bezüglich der Höhenfunktion $t: \R^3 \to \R$ (alle kritischen Punkte
von $t|_K$ sind nicht-entartet).

\begin{definition_de}[Kontsevich-Integral, analytische Form]\label{def:KI}
Für einen Knoten $K$ in allgemeiner Lage bezüglich der Höhenfunktion $t$
sei $z = x + iy$ für die ersten beiden Koordinaten.  Das
\emph{Kontsevich-Integral} ist:
\[
  Z(K) = \sum_{m=0}^\infty \frac{1}{(2\pi i)^m}
  \int_{t_{\min} < t_1 < \cdots < t_m < t_{\max}}
  \sum_{\substack{P = \{(z_j, z_j')\}_{j=1}^m \\ z_j(t_j), z_j'(t_j) \in K}}
  (-1)^{\downarrow P} D_P \; \bigwedge_{j=1}^m \frac{d z_j - d z_j'}{z_j - z_j'},
\]
wobei:
\begin{itemize}
  \item Die äußere Summe über $m = 0, 1, 2, \ldots$ läuft.
  \item Die innere Summe über alle Möglichkeiten läuft, $m$ horizontale
        Sehnen $(z_j(t_j), z_j'(t_j))$ mit verschiedenen Höhen
        $t_1 < \cdots < t_m$ zu wählen.
  \item $(-1)^{\downarrow P}$ gleich $(-1)$ hoch der Anzahl der Punkte in $P$
        ist, die auf einem abwärts orientierten Strang von $K$ liegen.
  \item $D_P$ das durch die Paarung $P$ bestimmte Sehnendiagramm ist.
  \item Die Form $\bigwedge_{j=1}^m \frac{dz_j - dz_j'}{z_j - z_j'}$ über
        das Simplex $t_{\min} < t_1 < \cdots < t_m < t_{\max}$ integriert wird.
\end{itemize}
\end{definition_de}

\begin{satz}[Kontsevich 1993, Invarianz]\label{satz:KI_invariant}
Das Kontsevich-Integral $Z(K) \in \hat{\mathcal{A}} = \prod_{n \geq 0} \mathcal{A}_n$
(die vervollständigte Sehnendiagramm-Algebra) ist eine Invariante
gerahmter Knoten.  Modulo der Gestellkorrektur ist es eine Invariante
nicht-gerahmter Knoten.
\end{satz}

\begin{proof}[Beweisidee]
\begin{enumerate}[label=(\arabic*)]
  \item $Z(K)$ ist invariant unter horizontalen Isotopien (Diffeomorphismen
        der Ebene bei jeder Höhe $t$): Dies folgt daraus, dass
        $\frac{dz - dz'}{z - z'}$ eine geschlossene $1$-Form ist.
  \item $Z(K)$ ist invariant unter den Reidemeister-Zügen II und III
        (aber \emph{nicht} Reidemeister I, der das Gestell um $1$ ändert):
        Dies erfordert eine sorgfältige Analyse in der Nähe kritischer Punkte,
        unter Verwendung des Assoziators.
  \item Die Gestellkorrektur verwendet das isolierte Sehnendiagramm $\theta$.
\end{enumerate}
Vollständige Details in Kontsevich~\cite{kontsevich1993} und
Bar-Natan~\cite{barnatan1995}.
\end{proof}

\begin{bemerkung}[Interpretation über die Magnus-Entwicklung]
Das Kontsevich-Integral kann auch als Holonomie einer flachen Verbindung
auf dem Konfigurationsraum $\mathrm{Konf}_n(\C)$ von $n$ Punkten in der
komplexen Ebene interpretiert werden.  Diese Verbindung ist die
\emph{Knizhnik-Zamolodchikov-Verbindung (KZ)}, die Werte in $\hat{\mathcal{A}}$
annimmt.
\end{bemerkung}

% -------------------------------------------------------
\section{Konvergenz und Regularisierung}\label{sec:konvergenz}
% -------------------------------------------------------

Das Integral in Definition~\ref{def:KI} hat logarithmische Divergenzen
in der Nähe kritischer Punkte von $K$.  Diese müssen regularisiert werden.

\begin{definition_de}[Beitrag kritischer Punkte]
In der Nähe eines lokalen Maximums (bzw.\ Minimums) von $K$ bei Höhe $t_0$
divergiert das Integral wie $\int \frac{dt}{t_0 - t} = -\log(t_0 - t)$.
Regularisiere durch:
\[
  Z(K) = \lim_{\varepsilon \to 0} \left[
    \text{(Integral über } t_{\min}+\varepsilon < \cdots < t_{\max}-\varepsilon)
    \cdot \exp\!\left(-\frac{1}{2} w(\varepsilon) \cdot \Phi\right)
  \right],
\]
wobei $w(\varepsilon)$ die ``Windung'' in der Nähe kritischer Punkte zählt
und $\Phi$ der Drinfeld-Assoziator ist.
\end{definition_de}

\begin{satz}[Le-Murakami 1996, Konvergenz]\label{satz:konvergenz}
Nach der obigen Regularisierung konvergiert das Kontsevich-Integral
und definiert ein wohldefiniertes Element von $\hat{\mathcal{A}}$ für
jeden Morse-Knoten $K$.
\end{satz}

% -------------------------------------------------------
\section{Der Drinfeld-Assoziator $\Phi_{\mathrm{KZ}}$}\label{sec:assoziator}
% -------------------------------------------------------

\begin{definition_de}[KZ-Gleichung]
Die \emph{Knizhnik-Zamolodchikov-Gleichung (KZ)} für eine Funktion
$G: \C\setminus\{0,1\} \to \hat{\mathcal{A}}$ ist:
\[
  \frac{dG}{dz} = \left(\frac{A}{z} + \frac{B}{z-1}\right) G,
\]
wobei $A, B$ formale Variablen (Erzeuger der Algebra $\hat{\mathcal{A}}$)
sind.  $A$ entspricht der Sehne bei $0$ und $B$ der Sehne bei $1$.
\end{definition_de}

\begin{satz}[Drinfeld 1990, Drinfeld-Assoziator]\label{satz:drinfeld}
Die KZ-Gleichung hat eine eindeutige Lösung $G_0(z)$, die asymptotisch
$z^A$ für $z \to 0^+$ ist (auf der reellen Achse), und eine eindeutige
Lösung $G_1(z)$, die asymptotisch $(1-z)^B$ für $z \to 1^-$ ist.
Der \emph{Drinfeld-Assoziator} (oder \emph{KZ-Assoziator}) ist:
\[
  \Phi_{\mathrm{KZ}} = G_0(z)^{-1} G_1(z) \in \hat{\mathcal{A}}.
\]
Er ist ein gruppenartiges Element in der vervollständigten Hopf-Algebra
$\hat{\mathcal{A}}$.
\end{satz}

\begin{proposition}[Eigenschaften von $\Phi_{\mathrm{KZ}}$]
\begin{enumerate}[label=(\roman*)]
  \item $\Phi_{\mathrm{KZ}}$ ist ein Assoziator: Er erfüllt die Fünfeck-Gleichung
        und die Sechseck-Gleichungen.
  \item $\Phi_{\mathrm{KZ}}$ ist gruppenartig: $\Phi_{\mathrm{KZ}} = \exp(\sum_n c_n [\text{Lie-Element vom Grad } n])$,
        wobei die Koeffizienten $c_n \in \Q$ Multiple-Zeta-Werte enthalten.
  \item Der Koeffizient des niedrigsten Grads von $\log \Phi_{\mathrm{KZ}}$
        in Grad $3$ ist $\frac{1}{24}[A, B]$ (verwandt mit $\zeta(3)$).
\end{enumerate}
\end{proposition}

\begin{bemerkung}[Assoziator und Multiple-Zeta-Werte]
Die Koeffizienten von $\Phi_{\mathrm{KZ}}$ in der Basis der Jacobi-Diagramme
sind (regularisierte) Multiple-Zeta-Werte (MZW):
$\zeta(n_1, \ldots, n_k) = \sum_{m_1 > \cdots > m_k > 0} m_1^{-n_1} \cdots m_k^{-n_k}$.
Dies ist eine der tiefen Verbindungen zwischen Knotentheorie und Zahlentheorie.
\end{bemerkung}

% -------------------------------------------------------
\section{Das farbige Jones-Polynom als Vassiliev-Invariante}\label{sec:jones}
% -------------------------------------------------------

\begin{definition_de}[Farbiges Jones-Polynom]
Für einen Knoten $K$ und eine ganze Zahl $N \geq 1$ ist das
\emph{$N$-farbige Jones-Polynom} $J_N(K; q)$ die Quanten-Invariante
zur $N$-dimensionalen irreduziblen Darstellung $V_N$ der Quantengruppe
$U_q(\mathfrak{sl}_2)$, normiert so dass
$J_N(\text{Unknoten}; q) = 1$.
\end{definition_de}

\begin{satz}[Vassiliev-Invarianten aus dem Jones-Polynom]
\label{satz:jones_vassiliev}
Das $N$-farbige Jones-Polynom $J_N(K;e^h)$, als Potenzreihe in $h$
entwickelt:
\[
  J_N(K; e^h) = \sum_{n=0}^\infty v_n^{(N)}(K) h^n,
\]
liefert Vassiliev-Invarianten $v_n^{(N)}$ vom Typ $\leq n$ für jedes $n$.
Das Gewichtssystem von $v_n^{(N)}$ ist $W_{\mathfrak{sl}_2, V_N}$.
\end{satz}

\begin{satz}[Kontsevich-Integral spezialisiert zum farbigen Jones-Polynom]
\label{satz:KI_jones}
Sei $W_{\mathfrak{sl}_2, V_N}: \hat{\mathcal{A}} \to \C[[h]]$ das
Lie-algebraische Gewichtssystem.  Dann gilt:
\[
  W_{\mathfrak{sl}_2, V_N}(Z(K)) = J_N(K; e^h).
\]
Das farbige Jones-Polynom ist also die Spezialisierung des Kontsevich-Integrals
unter dem Gewichtssystem $W_{\mathfrak{sl}_2, V_N}$.
\end{satz}

\begin{bemerkung}[Volumen-Vermutung]
Die \emph{Volumen-Vermutung} (Kashaev 1995, Murakami-Murakami 2001) besagt,
dass für einen hyperbolischen Knoten $K$:
\[
  \lim_{N \to \infty} \frac{2\pi}{N} \log|J_N(K; e^{2\pi i/N})| = \mathrm{vol}(S^3 \setminus K),
\]
wobei $\mathrm{vol}$ das hyperbolische Volumen bezeichnet.  Dies verbindet
die Vassiliev-/Kontsevich-Welt mit der hyperbolischen Geometrie.
\end{bemerkung}

% -------------------------------------------------------
\section{Universalität des Kontsevich-Integrals}\label{sec:universalitaet}
% -------------------------------------------------------

\begin{satz}[Kontsevich, Universalität]\label{satz:universalitaet}
Das Kontsevich-Integral $Z: \{\text{Knoten}\} \to \hat{\mathcal{A}}$ ist
die \emph{universelle Vassiliev-Invariante}: Jede Vassiliev-Invariante
faktoriert durch $Z$.  Genauer: Für jedes Gewichtssystem $W \in \hat{\mathcal{A}}^*$
ist die Komposition $W \circ Z: \{\text{Knoten}\} \to \mathbf{k}$ eine
Vassiliev-Invariante, und jede Vassiliev-Invariante entsteht auf diese Weise.
\end{satz}

\begin{proof}[Beweisidee]
Aus Satz~\ref{satz:vs_ws} folgt, dass Vassiliev-Invarianten durch
Gewichtssysteme klassifiziert werden.  Zu jedem Gewichtssystem $W$ sei
$v_W = W \circ Z$.  Dann ist $v_W$ eine Vassiliev-Invariante.  Umgekehrt
hat jedes $v \in \mathcal{V}_n$ ein Symbol $\bar{v} \in \Ac_n^*$, und
$v = \bar{v} \circ Z_n$ modulo Termen niedrigerer Ordnung (per Induktion
über $n$, unter Verwendung, dass $Z_n$ den Raum $\Ac_n$ über rationalen
Funktionen von $K$ aufspannt).
\end{proof}

\begin{korollar}
Das Kontsevich-Integral kodiert \emph{alle} Quantengruppen-Invarianten
von Knoten (alle farbigen Jones-Polynome, HOMFLY-PT-Polynom,
Kauffman-Polynom usw.) als Spezialisierungen.
\end{korollar}

% -------------------------------------------------------
\section{Die Vollständigkeitsvermutung}\label{sec:vollstaendigkeit}
% -------------------------------------------------------

\begin{vermutung}[Kontsevich-Integral trennt Knoten]\label{verm:vollstaendig}
Das Kontsevich-Integral $Z: \{\text{Knoten}\} / \text{Isotopie} \to \hat{\mathcal{A}}$
ist injektiv.  Das heißt:
\[
  Z(K_1) = Z(K_2) \implies K_1 \text{ und } K_2 \text{ sind isotop.}
\]
\end{vermutung}

\begin{bemerkung}
Dies würde bedeuten, dass Vassiliev-Invarianten ein \emph{vollständiges}
System von Knoten-Invarianten bilden: Sie unterscheiden gemeinsam alle Knoten.
Die Vermutung ist weitverbreitet geglaubt, aber vollständig offen.  Selbst
die schwächere Aussage, dass Vassiliev-Invarianten alle Knoten vom Unknoten
unterscheiden, ist unbekannt.
\end{bemerkung}

\begin{proposition}[Evidenz für Vermutung~\ref{verm:vollstaendig}]
\begin{enumerate}[label=(\roman*)]
  \item Das Kontsevich-Integral unterscheidet alle \emph{Primknoten} bis
        zu mindestens 10 Kreuzungen (rechnerisch verifiziert von
        Bar-Natan~\cite{barnatan1995}).
  \item $Z(K)$ unterscheidet Knoten von ihren Spiegelbildern, wenn $K$
        chiral ist (z.B.\ der Kleeblatt-Knoten).
  \item Es wurden keine zwei verschiedenen Knoten mit weniger als $10$
        Kreuzungen gefunden, die dasselbe Kontsevich-Integral haben.
\end{enumerate}
\end{proposition}

\begin{vermutung}[Vassiliev-Invarianten erkennen den Unknoten]\label{verm:unknoten}
Wenn alle Vassiliev-Invarianten von $K$ mit denen des Unknotens
übereinstimmen, dann ist $K$ der Unknoten.
\end{vermutung}

Diese schwächere Form von Vermutung~\ref{verm:vollstaendig} ist ebenfalls offen.

% -------------------------------------------------------
\section{Der Wheeling-Satz}\label{sec:wheeling}
% -------------------------------------------------------

Der Wheeling-Satz, bewiesen von Bar-Natan, Le und Thurston~\cite{barnatan2003},
ist eines der tiefsten Ergebnisse in der Theorie der Vassiliev-Invarianten.

\begin{definition_de}[Das Rad $\omega$ und die Wheeling-Abbildung]
In der Algebra $\mathcal{B}$ der Jacobi-Diagramme ist das \emph{Rad}
$\omega_n$ das Diagramm mit einer ``Nabe'' und $n$ Speichen.
Definiere:
\[
  \Omega = \sum_{n=0}^\infty \frac{b_{2n}}{n!} \omega_{2n}
  \in \hat{\mathcal{B}},
\]
wobei $b_{2n}$ die modifizierten Bernoulli-Zahlen sind
($\frac{t}{e^t-1} = \sum_n b_n t^n/n!$).
Die \emph{Wheeling-Abbildung} $\Phi_\Omega: \mathcal{B} \to \mathcal{A}$
ist durch das Einsetzen der Räder $\Omega$ an geeigneten Beinen definiert.
\end{definition_de}

\begin{satz}[Wheeling-Satz, Bar-Natan--Le--Thurston 2003]\label{satz:wheeling}
Die Wheeling-Abbildung
\[
  \Phi_\Omega: (\mathcal{B}, \cdot_{\text{disjunkt}}) \to (\mathcal{A}, \cdot_{\text{verbundene Summe}})
\]
ist ein Isomorphismus graduierter Algebren.  Dabei trägt $\mathcal{B}$
das disjunkte Vereinigungsprodukt und $\mathcal{A}$ das Verbundensummen-Produkt.
\end{satz}

\begin{proof}[Beweisidee]
Der Beweis zeigt, dass $\Phi_\Omega$ eine Algebra-Abbildung ist, unter
Verwendung der Kombinatorik der STU-Relation und der Eigenschaften der
Bernoulli-Zahlen.  Der Schlüsselschritt ist eine explizite Formel, die
die beiden Produkte über eine ``Wheeling''-Operation in Beziehung setzt.
\end{proof}

\begin{korollar}[Duflo-Isomorphismus als Spezialisierung]
Der Wheeling-Satz ist eine ``Kategorisierung'' des Duflo-Isomorphismus
in der Lie-Theorie.  Durch Anwenden des Gewichtssystems $W_{\mathfrak{g}}$
für eine Lie-Algebra $\mathfrak{g}$ auf den Wheeling-Satz erhält man
den Duflo-Isomorphismus
$\mathrm{Sym}(\mathfrak{g})^{\mathfrak{g}} \cong Z(U\mathfrak{g})$.
\end{korollar}

% -------------------------------------------------------
\section{Die Le-Murakami-Ohtsuki-Invariante}\label{sec:LMO}
% -------------------------------------------------------

Die Le-Murakami-Ohtsuki-Invariante (LMO) verallgemeinert das
Kontsevich-Integral auf $3$-Mannigfaltigkeiten.

\begin{definition_de}[Chirurgiedarstellung und LMO-Invariante]
Jede geschlossene orientierbare $3$-Mannigfaltigkeit $M$ kann durch
Chirurgie an einem Verschlingung $L \subset S^3$ erhalten werden.
Die \emph{LMO-Invariante} $Z_{\mathrm{LMO}}(M) \in \hat{\mathcal{A}}$
ist durch:
\[
  Z_{\mathrm{LMO}}(M) = \iota_n\left(Z(L) / \iota_n(Z(\bigcirc))^{\otimes n}\right)^{\otimes \sigma^\pm}
\]
definiert, wobei $n$ die Komponentenzahl von $L$, $\sigma^\pm$ die
Signatur der Verschlingungsmatrix, $\bigcirc$ der Unknoten und $\iota_n$
eine spezifische ``Entgestellungs-''Abbildung ist.
Genaue Definition in Le-Murakami-Ohtsuki~\cite{lmo1998}.
\end{definition_de}

\begin{satz}[Le-Murakami-Ohtsuki 1998]\label{satz:LMO}
Die LMO-Invariante $Z_{\mathrm{LMO}}(M)$ ist eine wohldefinierte Invariante
geschlossener orientierter $3$-Mannigfaltigkeiten.  Sie dominiert alle
Quanten-Invarianten von $3$-Mannigfaltigkeiten in dem Sinne, dass die
Reshetikhin-Turaev-Invarianten (für alle Quantengruppen) Spezialisierungen
von $Z_{\mathrm{LMO}}$ sind.
\end{satz}

\begin{bemerkung}
Die LMO-Invariante wird vermutet als vollständige Invariante rationaler
Homologie-Sphären zu sein (bis auf orientierungserhaltenden Diffeomorphismus),
aber dies ist offen.
\end{bemerkung}

% -------------------------------------------------------
\section{Bezug zur Chern-Simons-Feldtheorie}\label{sec:chern}
% -------------------------------------------------------

Witten~\cite{witten1989} schlug vor, dass alle Quanten-Invarianten von
Knoten und $3$-Mannigfaltigkeiten aus dem Pfadintegral stammen:
\[
  Z_{\mathrm{CS}}(M, K) = \int_{\mathcal{A}} e^{\frac{ik}{4\pi} \int_M \mathrm{Spur}(A \wedge dA + \tfrac{2}{3} A \wedge A \wedge A)} W_R(K) \,\mathcal{D}A,
\]
wobei $A$ ein Eichfeld (Zusammenhang) ist, $W_R(K) = \mathrm{Spur}_R \mathrm{Hol}_A(K)$
das Wilson-Schleifen-Observable in Darstellung $R$ und das Integral über
alle Zusammenhänge $\mathcal{A}$ modulo Eichinvarianz läuft.

\begin{satz}[Axiomatische Reformulierung, Reshetikhin-Turaev 1991]
Das Chern-Simons-Pfadintegral, streng gemacht durch die
Reshetikhin-Turaev-Konstruktion~\cite{reshetikhin1991} mit Quantengruppen
an Einheitswurzeln, liefert topologische Invarianten von $3$-Mannigfaltigkeiten
und Knoten darin.
\end{satz}

\begin{proposition}[Kontsevich-Integral als störungstheoretisches Chern-Simons]
Das Kontsevich-Integral $Z(K)$ entsteht als störungstheoretische Entwicklung
des Chern-Simons-Pfadintegrals um den trivialen Zusammenhang:
\[
  Z(K) = \sum_{n=0}^\infty \hbar^n \cdot [\text{Feynman-Diagramm-Beitrag der Ordnung } n].
\]
Die Feynman-Diagramme sind genau die Jacobi-Diagramme.
\end{proposition}

\begin{bemerkung}[Physikalische Interpretation]
Die $4T$-Relation in der Sehnendiagramm-Algebra entspricht der Jacobi-Identität
für die Lie-Klammer in $\mathfrak{g}$, die aus der Antisymmetrie der
Strukturkonstanten $f_{abc}$ entsteht.  Dies erklärt, warum Lie-algebraische
Gewichtssysteme natürlich sind.
\end{bemerkung}

% -------------------------------------------------------
\section{Offene Probleme}\label{sec:offen}
% -------------------------------------------------------

\begin{vermutung}[Vollständigkeit der Vassiliev-Invarianten, Birman-Lin 1993]
\label{verm:vollstaendigkeit2}
Vassiliev-Invarianten trennen gemeinsam alle Knoten.  Das heißt: Wenn
zwei Knoten $K_1, K_2$ für alle Vassiliev-Invarianten $v$ gleich sind
($v(K_1) = v(K_2)$), dann sind $K_1$ und $K_2$ isotop.
\end{vermutung}

\begin{vermutung}[Volumen-Vermutung, Kashaev-Murakami-Murakami]
\label{verm:volumen}
Für einen hyperbolischen Knoten $K$:
\[
  \lim_{N \to \infty} \frac{2\pi}{N} \log |J_N(K; e^{2\pi i/N})| = \mathrm{vol}(S^3 \setminus K).
\]
\end{vermutung}

\begin{vermutung}[Alle Gewichtssysteme sind Lie-algebraisch]
\label{verm:lie_alg}
Jedes Gewichtssystem stammt von einer metriseierten Lie-Algebra und
einer ihrer Darstellungen.
\end{vermutung}

\begin{bemerkung}[Bezug zur Kategorisierung]
Khovanov-Homologie~\cite{khovanov2000} ist eine Kategorisierung des
Jones-Polynoms: Sie liefert eine bigraduierte Homologietheorie
$\mathrm{Kh}^{i,j}(K)$ mit $J_2(K;q) = \sum_{i,j} (-1)^i q^j \dim \mathrm{Kh}^{i,j}(K)$.
Eine Kategorisierung des vollständigen Kontsevich-Integrals ist nicht bekannt ---
ein wichtiges offenes Problem.
\end{bemerkung}

\begin{thebibliography}{99}

\bibitem{vassiliev1990}
V.\ A.\ Vassiliev,
\textit{Cohomology of knot spaces},
in: Theory of Singularities and Its Applications (V.\ I.\ Arnold, Hrsg.),
Adv.\ Soviet Math.\ \textbf{1}, Amer.\ Math.\ Soc., 1990, S.\ 23--69.

\bibitem{kontsevich1993}
M.\ Kontsevich,
\textit{Vassiliev's knot invariants},
in: I.\ M.\ Gel'fand Seminar, Adv.\ Soviet Math.\ \textbf{16}, Teil 2,
Amer.\ Math.\ Soc., Providence, RI, 1993, S.\ 137--150.

\bibitem{barnatan1995}
D.\ Bar-Natan,
\textit{On the Vassiliev knot invariants},
Topology \textbf{34} (1995), 423--472.

\bibitem{jones1985}
V.\ F.\ R.\ Jones,
\textit{A polynomial invariant for knots via von Neumann algebras},
Bull.\ Amer.\ Math.\ Soc.\ \textbf{12} (1985), 103--111.

\bibitem{witten1989}
E.\ Witten,
\textit{Quantum field theory and the Jones polynomial},
Comm.\ Math.\ Phys.\ \textbf{121} (1989), 351--399.

\bibitem{drinfeld1990}
V.\ G.\ Drinfeld,
\textit{Quasi-Hopf algebras},
Leningrad Math.\ J.\ \textbf{1} (1990), 1419--1457.

\bibitem{birmanlin1993}
J.\ Birman, X.\ Lin,
\textit{Knot polynomials and Vassiliev's invariants},
Invent.\ Math.\ \textbf{111} (1993), 225--270.

\bibitem{barnatan2003}
D.\ Bar-Natan, S.\ Le, D.\ Thurston,
\textit{Two applications of elementary knot theory to Lie algebras and Vassiliev invariants},
Geom.\ Topol.\ \textbf{7} (2003), 1--31.

\bibitem{lmo1998}
T.\ Le, J.\ Murakami, T.\ Ohtsuki,
\textit{On a universal perturbative invariant of $3$-manifolds},
Topology \textbf{37} (1998), 539--574.

\bibitem{reshetikhin1991}
N.\ Reshetikhin, V.\ G.\ Turaev,
\textit{Invariants of $3$-manifolds via link polynomials and quantum groups},
Invent.\ Math.\ \textbf{103} (1991), 547--597.

\bibitem{khovanov2000}
M.\ Khovanov,
\textit{A categorification of the Jones polynomial},
Duke Math.\ J.\ \textbf{101} (2000), 359--426.

\bibitem{leturc1997}
T.\ Q.\ Le, J.\ Murakami,
\textit{The universal Vassiliev-Kontsevich invariant for framed oriented links},
Compositio Math.\ \textbf{102} (1996), 41--64.

\bibitem{barnatan1998}
D.\ Bar-Natan,
\textit{Non-associative tangles},
in: Geometric Topology (Proc.\ 1993 Georgia Internat.\ Topology Conf.),
AMS/IP Studies in Adv.\ Math.\ \textbf{2}, Amer.\ Math.\ Soc., 1997, S.\ 139--183.

\bibitem{chmutov2012}
S.\ Chmutov, S.\ Duzhin, J.\ Mostovoy,
\textit{Introduction to Vassiliev Knot Invariants},
Cambridge University Press, Cambridge, 2012.

\bibitem{ohtsuki2002}
T.\ Ohtsuki,
\textit{Quantum Invariants: A Study of Knots, 3-Manifolds, and Their Sets},
Series on Knots and Everything \textbf{29}, World Scientific, 2002.

\bibitem{murakami2001}
H.\ Murakami, J.\ Murakami,
\textit{The colored Jones polynomials and the simplicial volume of a knot},
Acta Math.\ \textbf{186} (2001), 85--104.

\end{thebibliography}

\end{document}
